% GENERATED FILE. Edit canonical lessons, then run npm run generate:books.
% canonical-source-hash: fnv1a64:1a0bfac60376d5f4
% canonical-lessons: MR-C339-nayak, MR-C339-khalnayak, MR-C339-patra, MR-C339-sadhu, MR-C339-savkar

\chapter{Heroes and Villains}
\label{ch:mr-heroes-and-villains}
\hlchaptermodality{339}

\begin{chapteropening}
\textbf{By the end of this chapter:} \emph{I can talk about the characters in a story.}

Complete the last lesson of chapter 339: I can talk about the characters in a story.
\end{chapteropening}

\section[\texorpdfstring{\mr{नायक} (nāyak)}{नायक (nāyak)}]{\mr{नायक} (nāyak) — a protagonist; the hero of a story}

\label{lesson:MR-C339-nayak}



\begin{quote}
Before the new one: say the Marathi for a witch, then the Marathi for a curse.
\end{quote}

\subsection*{You'll want to know: \mr{नायक}}
\textbf{\mr{नायक}} — \emph{nāyak} — \textquotedblleft{}a protagonist; the hero of a story\textquotedblright{}.

\begin{quote}
\mr{चित्रपटाचा} \mr{नायक} \mr{एक} \mr{साधा} \mr{शिक्षक} \mr{आहे}.
\end{quote}

\begin{quote}
\emph{citrapaṭācā nāyak ek sādhā śikṣak āhe.}
\end{quote}

\begin{quote}
The hero of the film is an ordinary teacher.
\end{quote}

Masculine.

A heroine is \mr{नायिका}.

\begin{cousinweb}
\textbf{\mr{नायक}} (\emph{nāyak}) is a tatsama, taken from Sanskrit \emph{nāyaka}, 'a guide, a leader; the hero of a drama' (Monier-Williams).
\end{cousinweb}

\begin{grammarlens}[title={What you've built}]
One more word for telling what happened.
\end{grammarlens}

\subsection*{Your turn}
\begin{itemize}
\raggedright
  \item \emph{Say it:} \emph{nāyak}
  \item \emph{Say it:} \emph{nāyak}, once more
  \item \emph{Say it:} \emph{citrapaṭācā nāyak ek sādhā śikṣak āhe.} — The hero of the film is an ordinary teacher.
  \item \emph{Recall:} say the Marathi for a witch, then the Marathi for a curse, then say \emph{nāyak} again
\end{itemize}

\begin{tcolorbox}[breakable,colback=teal!4,colframe=teal!35!black,title={Before you move on}]
What does \mr{नायक} mean? (\textquotedblleft{}A protagonist; the hero of a story\textquotedblright{}.) Say it once more.
\end{tcolorbox}
\section[\texorpdfstring{\mr{खलनायक} (khalnāyak)}{खलनायक (khalnāyak)}]{\mr{खलनायक} (khalnāyak) — a villain}

\label{lesson:MR-C339-khalnayak}



\begin{quote}
Before the new one: say the Marathi for a curse, then the Marathi for a protagonist.
\end{quote}

\subsection*{You'll want to know: \mr{खलनायक}}
\textbf{\mr{खलनायक}} — \emph{khalnāyak} — \textquotedblleft{}a villain\textquotedblright{}.

\begin{quote}
\mr{शेवटी} \mr{खलनायकाला} \mr{शिक्षा} \mr{झाली}.
\end{quote}

\begin{quote}
\emph{śevṭī khalnāyakālā śikṣā jhālī.}
\end{quote}

\begin{quote}
In the end the villain was punished.
\end{quote}

Masculine.

\mr{खल}, wicked, + \mr{नायक}, a lead.

\begin{cousinweb}
\textbf{\mr{खलनायक}} (\emph{khalnāyak}) is a modern compound built from Sanskrit \emph{khala}, 'a wicked man', and \emph{nāyaka}, 'a leader; the hero of a drama' (Monier-Williams).
\end{cousinweb}

\begin{grammarlens}[title={What you've built}]
One more word for telling what happened.
\end{grammarlens}

\subsection*{Your turn}
\begin{itemize}
\raggedright
  \item \emph{Say it:} \emph{khalnāyak}
  \item \emph{Say it:} \emph{khalnāyak}, once more
  \item \emph{Say it:} \emph{śevṭī khalnāyakālā śikṣā jhālī.} — In the end the villain was punished.
  \item \emph{Recall:} say the Marathi for a curse, then the Marathi for a protagonist, then say \emph{khalnāyak} again
\end{itemize}

\begin{tcolorbox}[breakable,colback=teal!4,colframe=teal!35!black,title={Before you move on}]
What does \mr{खलनायक} mean? (\textquotedblleft{}A villain\textquotedblright{}.) Say it once more.
\end{tcolorbox}
\section[\texorpdfstring{\mr{पात्र} (pātra)}{पात्र (pātra)}]{\mr{पात्र} (pātra) — a character}

\label{lesson:MR-C339-patra}



\begin{quote}
Before the new one: say the Marathi for a protagonist, then the Marathi for a villain.
\end{quote}

\subsection*{You'll want to know: \mr{पात्र}}
\textbf{\mr{पात्र}} — \emph{pātra} — \textquotedblleft{}a character\textquotedblright{}.

\begin{quote}
\mr{या} \mr{नाटकात} \mr{फक्त} \mr{तीन} \mr{पात्रं} \mr{आहेत}.
\end{quote}

\begin{quote}
\emph{yā nāṭkāt phakta tīn pātraṁ āhet.}
\end{quote}

\begin{quote}
There are only three characters in this play.
\end{quote}

Neuter.

Plural \mr{पात्रं}. It also means a vessel, and eligible.

\begin{cousinweb}
\textbf{\mr{पात्र}} (\emph{pātra}) is a tatsama, taken from Sanskrit \emph{pātra}, 'a vessel; a fit person; an actor' (Monier-Williams).
\end{cousinweb}

\begin{grammarlens}[title={What you've built}]
One more word for telling what happened.
\end{grammarlens}

\subsection*{Your turn}
\begin{itemize}
\raggedright
  \item \emph{Say it:} \emph{pātra}
  \item \emph{Say it:} \emph{pātra}, once more
  \item \emph{Say it:} \emph{yā nāṭkāt phakta tīn pātraṁ āhet.} — There are only three characters in this play.
  \item \emph{Recall:} say the Marathi for a protagonist, then the Marathi for a villain, then say \emph{pātra} again
\end{itemize}

\begin{tcolorbox}[breakable,colback=teal!4,colframe=teal!35!black,title={Before you move on}]
What does \mr{पात्र} mean? (\textquotedblleft{}A character\textquotedblright{}.) Say it once more.
\end{tcolorbox}
\section[\texorpdfstring{\mr{साधू} (sādhū)}{साधू (sādhū)}]{\mr{साधू} (sādhū) — a holy man; an ascetic}

\label{lesson:MR-C339-sadhu}



\begin{quote}
Before the new one: say the Marathi for a villain, then the Marathi for a character.
\end{quote}

\subsection*{You'll want to know: \mr{साधू}}
\textbf{\mr{साधू}} — \emph{sādhū} — \textquotedblleft{}a holy man; an ascetic\textquotedblright{}.

\begin{quote}
\mr{जंगलात} \mr{एक} \mr{म्हातारा} \mr{साधू} \mr{राहत} \mr{होता}.
\end{quote}

\begin{quote}
\emph{jaṅglāt ek mhātārā sādhū rāhat hotā.}
\end{quote}

\begin{quote}
An old holy man lived in the forest.
\end{quote}

Masculine.

Plural \mr{साधू}.

\begin{cousinweb}
\textbf{\mr{साधू}} (\emph{sādhū}) is a tatsama, taken from Sanskrit \emph{sādhu}, 'straight, good; a holy man' (Monier-Williams).
\end{cousinweb}

\begin{grammarlens}[title={What you've built}]
One more word for telling what happened.
\end{grammarlens}

\subsection*{Your turn}
\begin{itemize}
\raggedright
  \item \emph{Say it:} \emph{sādhū}
  \item \emph{Say it:} \emph{sādhū}, once more
  \item \emph{Say it:} \emph{jaṅglāt ek mhātārā sādhū rāhat hotā.} — An old holy man lived in the forest.
  \item \emph{Recall:} say the Marathi for a villain, then the Marathi for a character, then say \emph{sādhū} again
\end{itemize}

\begin{tcolorbox}[breakable,colback=teal!4,colframe=teal!35!black,title={Before you move on}]
What does \mr{साधू} mean? (\textquotedblleft{}A holy man; an ascetic\textquotedblright{}.) Say it once more.
\end{tcolorbox}
\section[\texorpdfstring{\mr{सावकार} (sāvkār)}{सावकार (sāvkār)}]{\mr{सावकार} (sāvkār) — a moneylender}

\label{lesson:MR-C339-savkar}



\begin{quote}
Before the new one: say the Marathi for a character, then the Marathi for a holy man.
\end{quote}

\subsection*{You'll want to know: \mr{सावकार}}
\textbf{\mr{सावकार}} — \emph{sāvkār} — \textquotedblleft{}a moneylender\textquotedblright{}.

\begin{quote}
\mr{गोष्टीतला} \mr{सावकार} \mr{गरीब} \mr{लोकांना} \mr{फार} \mr{त्रास} \mr{द्यायचा}.
\end{quote}

\begin{quote}
\emph{goṣṭītlā sāvkār garīb lokānnā phār trās dyāycā.}
\end{quote}

\begin{quote}
The moneylender in the story used to give poor people a lot of trouble.
\end{quote}

Masculine.

\mr{सावकाराकडून} \mr{कर्ज} \mr{घेणे}, to borrow from a moneylender.

\begin{cousinweb}
\textbf{\mr{सावकार}} (\emph{sāvkār}) is built on \emph{sāv}, 'a banker', which CDIAL lists under Sanskrit \emph{sādhu}, 'good' (CDIAL 13337), with the ending \emph{-kār}, 'one who does'.
\end{cousinweb}

\begin{grammarlens}[title={What you've built}]
One more word for telling what happened.
\end{grammarlens}

\subsection*{Your turn}
\begin{itemize}
\raggedright
  \item \emph{Say it:} \emph{sāvkār}
  \item \emph{Say it:} \emph{sāvkār}, once more
  \item \emph{Say it:} \emph{goṣṭītlā sāvkār garīb lokānnā phār trās dyāycā.} — The moneylender in the story used to give poor people a lot of trouble.
  \item \emph{Recall:} say the Marathi for a character, then the Marathi for a holy man, then say \emph{sāvkār} again
\end{itemize}

\begin{tcolorbox}[breakable,colback=teal!4,colframe=teal!35!black,title={Before you move on}]
What does \mr{सावकार} mean? (\textquotedblleft{}A moneylender\textquotedblright{}.) Say it once more.
\end{tcolorbox}
